%Begin


%%%%%%%%%%%%%%%%%%%%%%%%%%%%%%%%%%%%%%%%%%%%%%%%%%%%%%%%%%%%%%%
\chapter*{\S \space II. --- PRÉLIMINAIRES TOPOLOGIQUES~: SAUVAGERIE, DÉCOUPE}
\addcontentsline{toc}{section}{II. Préliminaires topologiques~: sauvagerie, découpe}
\section*{}

%%%%%%%%%%%%%%%%%%%%%%%%%%%%%%%%%%%%
\subsection*{1. Terminologie.}
\addcontentsline{toc}{subsection}{1. Terminologie}

{\bf 1.1}. --- 


\vskip .3cm
{\bf 1.2}. --- 

\vskip .3cm
{\bf 1.3}. --- 


\vskip .3cm
{\bf 1.4}. --- 

\vskip .3cm
{\bf 1.5}. --- 

\vskip .3cm
{\bf 1.6}. --- 



%%%%%%%%%%%%%%%%%%%%%%%%%%%%%%%%%%%%
\subsection*{2. Problèmes de sauvagerie.}
\addcontentsline{toc}{subsection}{2. Problèmes de sauvagerie}



\vskip .3cm
{\bf 2.1}. --- 


\vskip .3cm
{\bf 2.2}. --- Le lemme de segment.



\vskip .3cm
{\bf 2.3}. ---




%%%%%%%%%%%%%%%%%%%%%%%%%%%%%%%%%%%%
\subsection*{3. Isotopie à un plongement P.L.. Problèmes d'intersections.}
\addcontentsline{toc}{subsection}{3. Isotopie à un plongement P.L.. Problèmes d'intersections}

{\bf 3.1}. ---



\vskip .3cm
{\bf 3.2}. ---



%%%%%%%%%%%%%%%%%%%%%%%%%%%%%%%%%%%%
\subsection*{4. Ordres circulaires induits par un plongement de 1-complexe.}
\addcontentsline{toc}{subsection}{4. Ordres circulaires induits par un plongement de 1-complexe}

{\bf 4.1}. ---


\vskip .3cm
{\bf 4.2}. ---



\vskip .3cm
{\bf 4.3}. ---



%%%%%%%%%%%%%%%%%%%%%%%%%%%%%%%%%%%%
\subsection*{5. Découpe d'un espace topologique suivant un fermé.}
\addcontentsline{toc}{subsection}{5. Découpe d'un espace topologique suivant un fermé}

{\bf 5.1}. ---

\vskip .3cm
{\bf 5.2}. ---


\vskip .3cm
{\bf 5.3}. ---


\vskip .3cm
{\bf 5.4}. ---


\vskip .3cm
{\bf 5.5}. --- Découpe et espace quotient.



\vskip .3cm
{\bf 5.6}. --- Fonctorialité de la découpe.

La découpe ne faisant intervenir dans sa définition que des éléments topologiques (transportables par homéomorphisme) définit de fa\c{c}on évidente un foncteur de la catégorie des couples $(X, K)$ ($X$~: espace topologique, $K$~: fermé de $X$) munie des homéomorphismes de couples, dans elle-même. (On pourrait même prendre pour morphismes de $(X, K)$ dans $(X', K')$ toutes les applications continues $f: X \to X'$ telles que $f^{-1}(K') = K$).

Ce foncteur induit de fa\c{c}on toute aussi évidente un foncteur $\Sigma$ de la catégorie des couples $(X, K)$ ($X$~: surface à bord généralisé, $K$~: 1-complexe contenu dans $\operatorname{int} X$) munie des homéomorphismes de couples, dans la catégorie des surfaces à bord généralisé (munie des homéomorphismes), et ceci grâce au théorème {\bf 5.4.2.}.

Si $(X, K)$ et $(X', K')$ sont deux objets de la première catégorie, et $h$ un homéomorphisme de $(X, K)$ sur $(X', K')$ on en déduit fonctoriellement un homéomorphisme $\Sigma (h)$ de $\Sigma (X, K)$ sur $\Sigma (X', K')$ (homéomorphisme qui induit un homéomorphisme entre les bords généralisés de $\Sigma (X, K)$ et $\Sigma (X', K')$).

N.B.~: dans tout la suite, on gardera les notations $\Sigma (X, K)$ et $\Sigma (h)$.


% End
